% ===================================================================
%  CADRO N.º 10 — O mar. — O porto.
%  Traduction 2026 d'apres texto/io, controlee sur texto/fr, texto/es
%  et texto/pt. Consigne : voir texto/gl/10-cadro-01.tex.
%
%  Le francais decale sa numerotation d'un cran a partir de l'alinea
%  8. C'est une coquille de l'imprimeur francais ; le galicien suit
%  l'ido.
%
%  LE TABLEAU DE LA MER EST LE PLUS EMPRUNTE DE TOUTES LES COLONNES,
%  ET LE GALICIEN NE FAIT PAS EXCEPTION : « acoirazado », « torpedo »,
%  « cruceiro », « fragata », « dique », « hélice », « pilota »,
%  « semáforo », « draga » sont les memes mots qu'ailleurs. La marine
%  a vapeur est arrivee toute faite au XIXe siecle, avec ses noms.
%
%  ET LE PARTAGE SE FAIT ICI A L'INTERIEUR DE LA MER, ce qu'aucune
%  des cinq colonnes precedentes n'avait montre. Elles opposaient la
%  marine internationale aux mots de la TERRE — la falaise, la plage,
%  le rivage. Celle-ci l'oppose aux mots de la COTE, qui sont
%  galiciens jusqu'au dernier : « cantil » la falaise, « enseada » la
%  baie, « rompente » le brisant, « escollo » l'ecueil, « vagallada »
%  la lame, « peirao » la jetee, « cais » le quai. Le navire de
%  guerre s'apprend a l'ecole navale et prend le nom qu'on lui donne
%  la ; le brisant s'apprend en le regardant, et garde le sien.
% ===================================================================

%%K t10-apar-1 apar
\VUsaut{4.0mm}
\VUtitre{15.0pt}{60}{CADRO N.\textsuperscript{o} 10}
\VUsaut{3.0mm}
\VUpk{11.4pt}{40}{O mar. --- O porto.}
\VUsaut{3.0mm}

%%K t10-01-1 p
1. --- No verán, mentres uns van buscar a calma e a frescura na montaña, outros
prefiren ir á beira do mar. Iso foi o que fixemos o ano pasado, pois queremos moito
\VUgras{o Océano} \textsuperscript{(1)}.

%%K t10-02-1 p
2. --- Pola miña parte, sinto un gran pracer contemplando esa inmensa extensión de auga
que se estende ata o \VUgras{horizonte} \textsuperscript{(2)}, e vendo marchar para
unha longa \VUgras{travesía} os enormes \VUgras{vapores} \textsuperscript{(3)} nos que
embarcan os viaxeiros que van a América.

%%K t10-03-1 p
3. --- Por iso meu pai, que coñece os meus gustos, escolle sempre, para pasar as
vacacións, unha praia moi tranquila, situada preto dun dos nosos grandes
\VUgras{portos de guerra.}

%%K t10-03-2 p
Durante a nosa última estadía, a \VUgras{escuadra} viñera fondear detrás do enorme
\VUgras{dique} \textsuperscript{(4)} que pecha e protexe o \VUgras{porto militar}
\textsuperscript{(5)}, e conseguín admirar á miña vontade os distintos \VUgras{navíos
de guerra}: o pequeno \VUgras{submarino} \textsuperscript{(10)} que mergulla e
desaparece baixo as augas, e tamén o enorme \VUgras{acoirazado}
\textsuperscript{(9)} que, ata agora, só temía os ataques do \VUgras{torpedeiro}
\textsuperscript{(8)}.

%%K t10-tit-2 sub
\VUsaut{3.0mm}
\VUpk{10.2pt}{40}{Os navíos pequenos.}
\VUsaut{3.0mm}

%%K t10-04-1 p
4. --- Este xa sabía atopar moi \VUgras{habilmente} o punto feble da grosa coiraza
metálica para meter alí o perigoso \VUgras{torpedo}, cuxa explosión causa a perda do
navío; pero, aínda así, era menos temible ca o submarino, pois os contratorpedeiros
perségueno doadamente.

%%K t10-05-1 p
5. --- Puiden ver tamén os rápidos \VUgras{cruceiros} \textsuperscript{(6)}
acoirazados, case tan invulnerables coma o verdadeiro acoirazado, varios
\VUgras{transportes} \textsuperscript{(12)}, tres ou catro \VUgras{canoneiras}
\textsuperscript{(7)} e, finalmente, unha \VUgras{fragata} \textsuperscript{(11)}.
\VUgras{O espectáculo desa frota, cando manobra para levantar áncoras, é o máis
interesante que hai}.

%%K t10-06-1 p
6. --- Que diferenza de construción entre esas grandes moles de aceiro e os elegantes
\VUgras{trimastros} ou os graciosos \VUgras{veleiros} \textsuperscript{(13)} aínda
hoxe empregados na mariña mercante, que eu vía entrar no porto, a todo pano, cando
paseaba polo \VUgras{peirao} \textsuperscript{(14)}. É certo que estes perdían moitas
das súas vantaxes \VUgras{cando} o \VUgras{vento} era contrario. Tiñan que arriar
todas as velas para volver entrar na dársena, e un \VUgras{remolcador}
\textsuperscript{(16)} era preciso para os ir buscar e traelos, coma simples
\VUgras{gabarras} \textsuperscript{(17)}, ata o cais.

%%K t10-tit-3 sub
\VUsaut{3.0mm}
\VUpk{10.2pt}{40}{O porto de comercio.}
\VUsaut{3.0mm}

%%K t10-07-1 p
7. --- O que ten de interesante o porto do que falo é que non comprende só un porto
de guerra, feito a man con obras xigantescas, senón tamén un marabilloso \VUgras{porto
de comercio} situado na \VUgras{desembocadura} dun gran río.

%%K t10-08-1 p
8. --- A lingua de terra que separa as dúas dársenas está fortificada e defendida por
unha serie de \VUgras{baterías} e de \VUgras{baluartes} que avanzan sobre o mar. Cómpre
un permiso especial para pasear polas \VUgras{murallas} \textsuperscript{(15)}.
Conseguírao sen esforzo por medio do meu tío, que é enxeñeiro dos \VUgras{estaleiros}
\textsuperscript{(18)}, e fun visitar os navíos que se constrúen baixo amplos
alpendres.

%%K t10-09-1 p
9. --- Mesmo asistín á botadura dun acoirazado, e lembro o intre de emoción que sentiu
toda a multitude cando a mole enorme, abandonada a si mesma, comezou a esvarar sobre a
súa grade e veu flotar na auga do río.

%%K t10-10-1 p
10. --- Para ir aos estaleiros atravésase o \VUgras{campo de manobras}
\textsuperscript{(19)}, onde os soldados da guarnición veñen cada día facer o
exercicio, e pásase por unha \VUgras{pasarela} \textsuperscript{(20)} de ferro que
comunica a explanada do desfile co \VUgras{arsenal} \textsuperscript{(21)}.

%%K t10-tit-4 sub
\VUsaut{3.0mm}
\VUpk{10.2pt}{40}{O arsenal e os diques.}
\VUsaut{3.0mm}

%%K t10-11-1 p
11. --- Co meu tío visitei ese gran establecemento cheo de \VUgras{canóns}
\textsuperscript{(22)}, de \VUgras{balas} \textsuperscript{(23)} e de \VUgras{obuses}.
Vin tamén a \VUgras{fundición} \textsuperscript{(24)} na que se fabrican as
\VUgras{caldeiras} \textsuperscript{(25)} dos vapores.

%%K t10-12-1 p
12. --- Moi preto de alí están os \VUgras{diques} \textsuperscript{(26)} --- diques de
reparación --- e os moi curiosos \VUgras{diques flotantes} \textsuperscript{(27)} nos
que puiden ver de preto, nun navío en reparación, a \VUgras{hélice}
\textsuperscript{(28)}, a \VUgras{quilla} \textsuperscript{(29)} e o \VUgras{temón}
\textsuperscript{(30)} dun barco.

%%K t10-13-1 p
13. --- O porto de comercio é tan digno de estudo coma o porto de guerra, e pasei
moitas horas a papar moscas polos cais onde están amarrados os \VUgras{barcos de
rodas} \textsuperscript{(31)} que remontan o río, e mirando funcionar as \VUgras{guindastres
de mastrear} \textsuperscript{(32)}, que erguen coma plumas cargas enormes.

%%K t10-tit-5 sub
\VUsaut{3.0mm}
\VUpk{10.2pt}{40}{O práctico.}
\VUsaut{3.0mm}

%%K t10-14-1 p
14. --- Admirei os \VUgras{transatlánticos} \textsuperscript{(34)} que saían da rada,
coas súas \VUgras{bandeiras} \textsuperscript{(35)} ao vento, guiados por un hábil
\VUgras{práctico} \textsuperscript{(36)} que sabía marabillosamente seguir a
\VUgras{canle} marcada por \VUgras{boias} \textsuperscript{(40)} e \VUgras{balizas},
evitando as \VUgras{gabarras} \textsuperscript{(41)} que se dirixen a custo coa súa
única vela cadrada. Distinguín na \VUgras{proa} \textsuperscript{(37)} --- a
amura --- os \VUgras{camarotes} \textsuperscript{(39)} de segunda clase, e na
\VUgras{popa} \textsuperscript{(38)} --- o coroamento --- os salóns para os
pasaxeiros de primeira clase.

%%K t10-15-1 p
15. --- Ás veces asistía tamén á carga dun \VUgras{cargueiro} \textsuperscript{(42)}
no que se amoreaban as mercadorías do \VUgras{almacén} \textsuperscript{(43)} e da
\VUgras{estación marítima} \textsuperscript{(44)}, traídas polos moitos trens da
\VUgras{liña dos cais} \textsuperscript{(45)}.

%%K t10-tit-6 sub
\VUsaut{3.0mm}
\VUpk{10.2pt}{40}{O pracer de remar.}
\VUsaut{3.0mm}

%%K t10-16-1 p
16. --- Moitas veces andei en bote na \VUgras{dársena de flotación}
\textsuperscript{(53)}, na que veñen refuxiarse as \VUgras{embarcacións}
\textsuperscript{(46)}, e era unha ledicia para min manexar o \VUgras{remo}
\textsuperscript{(47)}. Subín á \VUgras{cuberta} \textsuperscript{(48)} dunha
\VUgras{barca de vela} \textsuperscript{(49)}, gabeei, coas \VUgras{cordas}
\textsuperscript{(52)}, polo \VUgras{mastro} \textsuperscript{(51)} ata as
\VUgras{vergas} \textsuperscript{(50)}, e despois seguín o meu paseo \VUgras{en iola}
\textsuperscript{(54)} entre as pesadas \VUgras{barcas de pesca}
\textsuperscript{(55)}, onde secan as \VUgras{redes} \textsuperscript{(56)}.

%%K t10-17-1 p
17. --- No \VUgras{cais} \textsuperscript{(58)} desfilaban diante de min as máis
diversas \VUgras{mercadorías} \textsuperscript{(59)}, metidas en \VUgras{caixas}
\textsuperscript{(60)} ou en \VUgras{fardos} \textsuperscript{(61)}. Formaban a
\VUgras{carga} \textsuperscript{(63)} das \VUgras{gabarras}, e uns
\VUgras{guindastres} potentes \textsuperscript{(62)} colléronas e depositáronas en
\VUgras{camións} \textsuperscript{(64)} mentres os \VUgras{transportistas} facían a
súa declaración e pagaban os dereitos na \VUgras{aduana} \textsuperscript{(65)}.

%%K t10-18-1 p
18. --- Finalmente, cando cheguei á \VUgras{ponte xiratoria} \textsuperscript{(66)} ou
á \VUgras{esclusa} \textsuperscript{(69)}, pasando diante da \VUgras{draga}
\textsuperscript{(68)} que limpa a \VUgras{canle} \textsuperscript{(67)}, desembarquei
no peirao preto do \VUgras{semáforo} \textsuperscript{(70)}.

%%K t10-19-1 p
19. --- É do outro lado dese peirao onde veñen atracar as \VUgras{chalupas}
\textsuperscript{(71)} que traen a terra os oficiais da escuadra. É un espectáculo
admirable ver os mariñeiros erguer a compás os seus remos, mentres a barca, levada pola
\VUgras{onda} \textsuperscript{(72)}, atraca sen esforzo na escaleira do
desembarcadoiro. Neste lugar é onde mellor se pode observar a \VUgras{marea} e o
movemento do \VUgras{fluxo} e do \VUgras{refluxo} sobre a \VUgras{praia}
\textsuperscript{(73)}.

%%K t10-tit-7 sub
\VUsaut{3.0mm}
\VUpk{10.2pt}{40}{O casino.}
\VUsaut{3.0mm}

%%K t10-20-1 p
20. --- Para volver á casa collín o \VUgras{tranvía eléctrico}
\textsuperscript{(74)}, cuxa \VUgras{parada} \textsuperscript{(75)} está preto do
\VUgras{poste} posto diante do círculo dos oficiais.

%%K t10-20-2 p
Desde o \VUgras{balcón} \textsuperscript{(76)} do círculo, adornado con
\VUgras{áncoras} \textsuperscript{(77)}, canóns e balas, vin moitas veces unha
espléndida xuntanza de oficiais de mariña: \VUgras{tenentes de navío}
\textsuperscript{(78)}, coa súa espada, \VUgras{capitáns de artillaría}
\textsuperscript{(79)} e de \VUgras{infantaría} \textsuperscript{(80)} con
\VUgras{sabre}. Detrás deles había un \VUgras{mariñeiro} \textsuperscript{(81)} que
lles traía un \VUgras{catalexo} cando querían distinguir o que sucedía ao lonxe.

%%K t10-21-1 p
21. --- Vin tamén novos \VUgras{alféreces} \textsuperscript{(82)} recibindo as
instrucións do seu \VUgras{comandante} \textsuperscript{(83)} ou do
\VUgras{almirante} \textsuperscript{(84)}, mentres un \VUgras{contramestre}
\textsuperscript{(85)} agardaba para as levar ao navío.

%%K t10-22-1 p
22. --- Desde ese balcón, a vista é espléndida: domínase toda a praia coas súas
\VUgras{tendas} \textsuperscript{(86)} e \VUgras{casetas} \textsuperscript{(87)}, e
vense, á hora do baño, os \VUgras{bañistas} \textsuperscript{(88)} que brincan ledos
diante do \VUgras{casino} \textsuperscript{(89)}.

%%K t10-23-1 p
23. --- Na punta da praia, a costa forma unha pequena \VUgras{península}
\textsuperscript{(91)} \VUgras{rochosa}, onde vén romper a \VUgras{vagallada}
\textsuperscript{(92)} e sobre a que está construído un \VUgras{faro}
\textsuperscript{(93)}, que indica de noite o camiño aos navegantes. Esa península só
comunica coa terra por un \VUgras{istmo} moi estreito \textsuperscript{(90)}.

%%K t10-tit-8 sub
\VUsaut{3.0mm}
\VUpk{10.2pt}{40}{Vista xeral da costa.}
\VUsaut{3.0mm}

%%K t10-24-1 p
24. --- O \VUgras{cantil} \textsuperscript{(94)} esténdese, fendido polo mar, que
debuxa nel caprichosamente unha serie de \VUgras{enseadas} \textsuperscript{(95)} e de
\VUgras{promontorios} (cabos), que o fan moi pintoresco.

%%K t10-24-2 p
Sobre o \VUgras{chan alto} \textsuperscript{(96)}, que domina o \VUgras{estreito}
\textsuperscript{(98)}, hai un \VUgras{forte} \textsuperscript{(97)} moi importante e
algunhas «villas».

%%K t10-25-1 p
25. --- Ese estreito separa do continente unha \VUgras{illa} \textsuperscript{(99)} moi
perigosa, cinguida de \VUgras{escollos} \textsuperscript{(100)} e de
\VUgras{rompentes}, onde barcas e mesmo grandes navíos encallan ás veces. A pesar da
prontitude dos socorros e da chegada rápida dos \VUgras{botes de salvamento}, eses
\VUgras{naufraxios} son terroríficos e, durante as \VUgras{tempestades} do equinoccio,
varios navíos pereceron con xente e mercadorías.

%%K t10-25-2 p
A pesar desa veciñanza, \VUgras{o porto} é moi frecuentado polos mariñeiros.
